\section{Trabalhos relacionados}

A qualidade dos sinais eletromiográficos de superfície (sEMG) é fator determinante para o sucesso de sistemas de reconhecimento de gestos e controle de próteses. A literatura recente tem concentrado esforços tanto no \textit{hardware} de aquisição, multiplicando canais e integrando filtros embarcados, quanto no pré‑processamento digital, etapa essencial para reduzir ruídos e artefatos antes da classificação.
A Tabela~\ref{tab:aquisicao_semg} sumariza os principais trabalhos que abordam a aquisição e a filtragem de sinais sEMG, evidenciando as soluções adotadas em diferentes épocas e contextos.

\begin{table}[htpb]
    \centering
    \small
    \setlength{\extrarowheight}{3pt}
    \caption{Trabalhos relacionados focados na aquisição e filtragem de sinais sEMG.}
    \label{tab:aquisicao_semg}
    \begin{tabularx}{\textwidth}{@{} 
        >{\RaggedRight\arraybackslash}X 
        >{\RaggedRight\arraybackslash}X 
        >{\RaggedRight\arraybackslash}X 
        @{}}
        \toprule
        \rowcolor{gray!20} \textbf{Título do Artigo (Ano)} & \textbf{Sensor (Dispositivo, Quantidade e Canais)} & \textbf{Filtros Usados (Pré-processamento)} \\
        \midrule
        \textit{Processing and Analysis of EMG Signals from MYO Armband for Upper-Limb Prosthesis Control} \cite{ProcessingAndAnalysisOfEMGSignalsFromMYO} & 1 dispositivo (Myo Armband). 8 canais de eletrodos a seco posicionados ao redor do antebraço. & Filtro Passa-banda de 20--90 Hz e Filtro \textit{Notch} de 50 Hz. Retificação e suavização (RMS). \\
        \addlinespace
        \textit{High-Performance Surface Electromyography Armband Design for Gesture Recognition} \cite{Zhang2023-gn} & 1 bracelete customizado ($\alpha$ Armband). 16 canais com conversores ADC de 16-bit. & Largura de banda ajustável de 0.1 Hz a 20 kHz. Filtragem no hardware para capturar altas frequências. \\
        \addlinespace
        \textit{Medium density EMG armband for gesture recognition} \cite{Aghchehli2025-te} & 1 bracelete de “Média Densidade” customizado. 21 canais usando eletrodos digitais de contato por mola. & Filtro de Interferência Eletromagnética (EMI) integrado no hardware. IA minimiza necessidade de filtros digitais. \\
        \addlinespace
        \textit{Machine Learning Approaches for EMG Signal Classification to Predict Hand Gesture} \cite{Rendon_Caballero2026-ar} & Eletrodos de superfície padrão. Comparação entre matrizes de 10, 7 e 4 canais. & Filtro \textit{Notch} de 50 Hz e Filtro Passa-banda de 50--500 Hz. Segmentação em janelas de tempo curtas. \\
        \addlinespace
        \textit{putEMG—A Surface Electromyography Hand Gesture Recognition Dataset}\cite{Kaczmarek2019-mw}  & Matriz de eletrodos sEMG no antebraço coletando dados em altíssima densidade via 24 canais. & Filtros digitais passa-baixo e passa-alto nos dados brutos para estabelecer o \textit{baseline} de classificação. \\
        \bottomrule
    \end{tabularx}
\end{table}

A análise conjunta dos trabalhos revela três tendências claras:
\begin{enumerate}
    \item \textbf{Densificação dos canais}: Os sistemas evoluíram de 8 canais (Myo Armband) para arranjos com 21 ou 24 canais, buscando maior resolução espacial do sinal.
    \item \textbf{Pré‑processamento \textit{on‑chip}}: A filtragem analógica ou digital embarcada (EMI, banda ajustável) vem ganhando espaço, reduzindo a carga sobre o software de classificação.
    \item \textbf{Par mínimo de filtros digitais}: A combinação de um filtro passa‑banda (tipicamente 20–500 Hz) com um filtro \textit{notch} em 50 Hz (ou 60 Hz) persiste como prática padrão, mesmo quando métodos de aprendizado de máquina são empregados na etapa seguinte.
\end{enumerate}
Essas constatações orientam a escolha do \textit{hardware} e da cadeia de filtros adotados no presente trabalho, que busca equilibrar densidade de canais, qualidade do sinal e complexidade computacional.